\clearpage
\section*{术语表}
\addcontentsline{toc}{section}{术语表}

下表列出本译本中反复出现的关键术语，按首次出现的大致顺序排列。凡原文给出缩写者，缩写一并列出；正文中首次出现处已给出“中文（English）”形式，其后统一使用中文或该缩写。人名、群名与定理名保留拉丁拼写，不列入本表。

\begin{center}
\small
\begin{longtable}{@{}>{\raggedright\arraybackslash}p{0.24\linewidth}>{\raggedright\arraybackslash}p{0.44\linewidth}>{\raggedright\arraybackslash}p{0.20\linewidth}@{}}
\toprule
中文 & English（缩写） & 首次出现 \\
\midrule
\endfirsthead
\toprule
中文 & English（缩写） & 首次出现 \\
\midrule
\endhead
\bottomrule
\endfoot
自守层 & automorphic sheaf & §1 \\
Hecke 本征层 & Hecke eigensheaf & §1 \\
本征值／本征空间 & eigenvalue / eigenspace & §1.1 \\
谱定理 & spectral theorem & §1.1 \\
投影值测度 & projection-valued measure & §1.1 \\
泛函演算 & functional calculus & §1.1 \\
直接积分 & direct integral & §1.1 \\
梅林变换 & Mellin transform & §1.2 \\
傅里叶--Malgrange 变换 & Fourier--Malgrange transform & §1.2 \\
范畴化 & categorification & §1.2 \\
朗兰兹对偶群 & Langlands dual group, $\Gd$ & §1.4 \\
朗兰兹参数 & Langlands parameter & §1.4 \\
Hecke 算子 & Hecke operator & §1.4 \\
尖点 & cusp & §1.4 \\
局部 Langlands 对应 & local Langlands correspondence, LLC & §1.4 \\
整体域／局部域 & global field / local field & §1.4 \\
表示 & representation & §2.1 \\
层 & sheaf & §2.1 \\
导出范畴 & derived category & §2.1 \\
六函子形式 & six-functor formalism & §2.1 \\
拟凝聚层 & quasi-coherent sheaf, $\QCoh$ & §2.1 \\
归纳凝聚层 & ind-coherent sheaf, $\IndCoh$ & §2.1 \\
模栈 & moduli stack & §2.2 \\
$G$-丛模栈 & moduli stack of $G$-bundles, $\Bun_{G}$ & §2.2 \\
半稳定丛／不稳定丛 & semistable / unstable bundle & §2.2.2 \\
雅可比簇 & Jacobian variety, $\mathrm{Jac}(C)$ & §2.2.1 \\
Picard 簇 & Picard variety & §2.2.1 \\
亏格 & genus, $g$ & §2.2.1 \\
仿射 Grassmannian & affine Grassmannian, $\Gr_{G}$ & §2.2.2 \\
阿代尔描述 & adelic description & §2.2.2 \\
Weil 单值化 & Weil uniformization & §2.2.2 \\
除子 & divisor & §2.2.3 \\
局部系统 & local system & §2.3 \\
Betti 版本／de Rham 版本 & Betti / de~Rham version & §2.3 \\
特征簇 & character variety & §2.3 \\
平坦联络 & flat connection & §2.3.1 \\
类域论 & class field theory & §2.4 \\
Abel--Jacobi 映射 & Abel--Jacobi map & §2.4 \\
D-模 & D-module & §2.4.1 \\
de Rham 共形场论 & de~Rham chiral CFT & §2.4.1 \\
手征代数 & chiral algebra & §2.4.1 \\
Fourier--Mukai 变换 & Fourier--Mukai transform & §2.4.1 \\
摩天楼层 & skyscraper sheaf & §2.4.1 \\
白光／Whittaker 层 & white field, $\mathrm{White}_{G}$ & §2.5 \\
Hecke 范畴 & Hecke category & §3.1 \\
球 Hecke 范畴 & spherical Hecke category, $\Sph_{G}$ & §3.1 \\
't Hooft 线算子 & 't~Hooft line operator & §3.1 \\
算子乘积展开 & operator product expansion, OPE & §3.2 \\
因子化代数 & factorization algebra & §3.2 \\
因子化同调 & factorization homology & §3.2 \\
小圆盘算子代数 & little discs operad & §3.2 \\
对称幺半范畴 & symmetric monoidal category & §3.2 \\
几何 Satake 等价 & geometric Satake equivalence & §3.3 \\
Tannaka 重构 & Tannaka reconstruction & §3.4 \\
谱作用 & spectral action & §3.5 \\
希格斯丛 & Higgs bundle, $\Higgs_{G}$ & §4.1 \\
Hitchin 映射／纤维化 & Hitchin map / fibration, $\Hitch$ & §4.1 \\
谱曲线 & spectral curve & §4.1 \\
Hitchin 截面 & Hitchin section & §4.1 \\
幂零锥 & nilpotent cone, $\Nilp$ & §4.1 \\
辛几何 & symplectic geometry & §4.1 \\
SYZ 镜像对称 & SYZ mirror symmetry & §4.1 \\
Dolbeault 几何朗兰兹 & Dolbeault geometric Langlands & §4.1 \\
微局部 & micro-local & §4.2 \\
奇异支撑 & singular support & §4.2 \\
阿贝尔化 & abelianization & §4.2 \\
交换微分算子 & commuting differential operators & §4.2 \\
非阿贝尔 Hodge 理论 & nonabelian Hodge theory & §4.2 \\
紧支撑 & compact support & §5.1 \\
自对偶 & self-dual & §5.1 \\
导出代数几何 & derived algebraic geometry & §5.1 \\
étale 版本 & étale version & §5.1.1 \\
周期 & period & §5.2 \\
相对 Langlands 纲领 & relative Langlands program & §5.2 \\
玉河数 & Tamagawa number & §5.2 \\
$\Theta$-级数 & theta series & §5.2 \\
Eisenstein 级数 & Eisenstein series & §5.3 \\
连续谱 & continuous spectrum & §5.3 \\
抛物子群／Levi 子群 & parabolic subgroup / Levi subgroup & §5.3 \\
Whittaker 周期 & Whittaker period & §5.4 \\
Nahm 极点 & Nahm pole & §5.4 \\
$q$-展开 & $q$-expansion & §5.4 \\
Kac--Moody 代数 & Kac--Moody algebra & §5.5 \\
临界水平 & critical level & §5.5 \\
Beilinson--Bernstein 局部化 & Beilinson--Bernstein localization & §5.5 \\
oper & oper & §5.5 \\
分歧 & ramification & §6.0.1 \\
温和分歧 & tame ramification & §6.0.1 \\
量子群 & quantum group, $U_{q}\mathfrak{g}$ & §6.0.2 \\
正特征 & positive characteristic & §6.0.2 \\
刚性局部系统 & rigid local system & §6.0.3 \\
反 Galois 问题 & inverse Galois problem & §6.0.3 \\
可积系统 & integrable system & §6.0.3 \\
量子几何朗兰兹 & quantum geometric Langlands & §6.0.4 \\
基本局部等价 & fundamental local equivalence & §6.0.4 \\
Kazhdan--Lusztig 等价 & Kazhdan--Lusztig equivalence & §6.0.4 \\
Fargues--Fontaine 曲线 & Fargues--Fontaine curve & §6.1.1 \\
算术拓扑 & arithmetic topology & §6.1.1 \\
twistor 空间 & twistor space, $\PP^{1}$ & §6.1.1 \\
$\ell$-独立性 & $\ell$-independence & §6.1.1 \\
函数域 & function field & §6.1.2 \\
Weil 的罗塞塔石碑 & Weil's Rosetta Stone & §6.1.2 \\
函数--层对应 & function-sheaf correspondence & §6.1.2 \\
基本引理 & fundamental lemma & §6.1.2 \\
Ramanujan 猜想 & Ramanujan conjecture & §6.1.2 \\
Montonen--Olive 对偶 & Montonen--Olive duality & §6.2 \\
超对称 & supersymmetry & §6.2 \\
拓扑扭曲 & topological twist & §6.2 \\
拓扑量子场论 & topological quantum field theory, TQFT & §6.2 \\
边界条件 & boundary condition & §6.2 \\
BPS 膜 & BPS brane & §6.2 \\
Hodge 模／twistor 模 & Hodge module / twistor module & §6.2 \\
motive & motive & §6.2 \\
\end{longtable}
\end{center}

